\answer{ex:18:1}{(a)~$0.9949\le w\le1.0049$，即 $|1-w|\lesssim0.005$。(b)~$K\ge400$ 个突触。}
\answer{ex:18:2}{$(\omega_R\tau_w)^2=\sqrt{\alpha(\alpha+2+2\varrho)}/\varrho-1$，$\varrho=\tau_m/\tau_w$；$f_R\approx11.1$~Hz；$|Z(\omega_R)|/|Z(0)|=4.6$。}
\answer{ex:18:3}{(a)~$f=1/\bigl(\tau\ln\frac{\mu}{\mu-\theta}\bigr)$；要得到 $1$~Hz，需要 $\mu/\theta-1\approx3.7\cdot10^{-44}$。(b)~$T=\pi\tau/\sqrt\mu$，$f\approx3.2$~Hz：$f\propto\sqrt\mu$。}
\answer{ex:18:4}{积分器在 $f\lesssim17.7$~Hz时工作（低通滤波器），谐振器只在 $5.5$--$22$~Hz频带内工作（带通滤波器）。}
\answer{ex:18:5}{(a)~$T_{\min}=\frac{\tau}{w-1}\ln\frac{0.5}{0.3}\approx10.2\ms$。(b)~$B>w-1-0.2=0.3$。}
\answer{ex:18:6}{调节的时间常数为 $\tau/(\eta\langle r^2\rangle)$，$\eta=\tau\ln10/60\,\text{s}\approx3.8\cdot10^{-3}$；$I\neq0$ 时权重会偏移，需要从 $e$ 中减去指令的副本 $I/\tau$。}
\answer{ex:18:7}{开放问题。谐振器在 $11$~Hz时仅胜出 $1.35$ 倍，在 $1$~Hz时则落后 $17$ 倍；重构的主要收益在于借助阈值选择频带（习题~\ref{ex:18:4}）。}
